%--- Part 4: Interpretable Modeling (8 min) ------------------------------------

\section{Interpretable Modeling}

\begin{frame}{Analysis Pipeline}
  \vfill\pipelinehighlight{2}\vfill
\end{frame}

%--- Slide: Why interpretable models? ------------------------------------------
\begin{frame}{Why Interpretable Models?}
  \begin{itemize}[<+->]
    \item We have per-trajectory \texttt{ml\_ade} — but raw error values
          alone do not tell us \emph{why} some trajectories are hard
    \item A black-box model can predict ADE well but gives no insight
          into which features drive it or in which direction
    \item \textbf{Interpretable models} let us directly read off feature
          effects: how does speed, heading change, or scene density
          affect prediction error?
    \item This turns ADE into actionable understanding of model
          behaviour
  \end{itemize}
\end{frame}

%--- Slide: GAM and XGBoost ---------------------------------------------------
\begin{frame}{GAM and XGBoost}
  We tested two interpretable modeling approaches --- either can be used
  for the analysis.
  \vspace{0.5em}
  \begin{columns}[t]
    \begin{column}{0.48\textwidth}
      \textbf{GAM (Generalized Additive Model)}
      \visible<2->{\begin{itemize}[<+->]
        \item Models target as a sum of smooth functions: $f(x) = \sum_j s_j(x_j)$
        \item One spline per feature — no interaction terms
        \item p-values for significance testing
        \item Directly readable smooth effects
      \end{itemize}}
    \end{column}
    \begin{column}{0.48\textwidth}
      \textbf{XGBoost}
      \visible<3->{\begin{itemize}[<+->]
        \item Gradient-boosted decision trees
        \item Flexibly captures non-linearities
              and feature interactions
        \item Post-hoc interpretation via SHAP values
      \end{itemize}}
    \end{column}
  \end{columns}
\end{frame}

%--- Slide: Additive effects & SHAP -------------------------------------------
\begin{frame}{Interpreting the Models: Additive Effects \& SHAP Values}
  \begin{columns}[t]
    \begin{column}{0.48\textwidth}
      \textbf{GAM — Additive Effects}
      \begin{itemize}[<+->]
        \item The model is a sum of smooth terms:
              $\hat{y} = \sum_j \hat{s}_j(x_j)$
        \item Each $\hat{s}_j$ directly shows how feature $j$
              affects \texttt{ml\_ade} — no post-hoc step needed
        \item Confidence bands from the spline fit
      \end{itemize}
    \end{column}
    \begin{column}{0.48\textwidth}
      \textbf{XGBoost — SHAP Values}
      \begin{itemize}[<+->]
        \item Each prediction is decomposed into per-feature
              contributions post-hoc
        \item \emph{Local:} why was this trajectory's ADE high?
        \item \emph{Global:} mean $|$SHAP$|$ gives feature
              importance and direction
      \end{itemize}
    \end{column}
  \end{columns}
  \vspace{0.5em}
\end{frame}

%--- Slide 12 ------------------------------------------------------------------
\begin{frame}{Validating the Tuning Procedure — Nested Resampling}
  \begin{itemize}
    \item \textbf{Problem:} hyperparameter tuning risks overfitting to a
          particular data split
    \item \textbf{Solution:} nested cross-validation
    \begin{itemize}
      \item Outer folds: evaluate generalization of the full tuning procedure
      \item Inner folds: tune hyperparameters on each outer-fold training set via cross-validation
    \end{itemize}
    \item OOF evaluation: inspect out-of-fold predictions for instability,
          outlier folds, bias
    \item \textbf{Conclusion:} tuning generalises stably $\to$ proceed to
          full-data tuning + refit
  \end{itemize}
  \vspace{0.4em}
  \footnotesize Final models are refit on \emph{all} data with optimal
  hyperparameters — no held-out test set, because the goal is understanding,
  not unbiased prediction.
\end{frame}

%--- Slide: GAM Variant Selection ----------------------------------------------
\begin{frame}{GAM Variant Selection}
  \begin{columns}[t]
    \begin{column}{0.44\textwidth}
      \textbf{Three candidates:}
      \begin{itemize}[<+->]
        \item \textbf{LinearGAM} — fitted on raw \texttt{ml\_ade}
        \item \textbf{LinearGAM (log)} — fitted on
              $\log(\texttt{ml\_ade})$, evaluated on original scale
        \item \textbf{GammaGAM} — Gamma distribution with log link,
              suited for positive right-skewed targets
      \end{itemize}
      \vspace{0.5em}
      \onslide<2->{Selection criterion: lowest mean outer-fold RMSE
      (original scale).}
    \end{column}
    \begin{column}{0.53\textwidth}
      \onslide<2->{
      \renewcommand{\arraystretch}{1.3}
      \begin{tabular}{lccc}
        \toprule
        \textbf{Variant} & \textbf{RMSE} & \textbf{R$^2$} & \textbf{MAE} \\
        \midrule
        \textcolor{lmugreen}{\textbf{LinearGAM}} & \textcolor{lmugreen}{\textbf{0.440}} & \textcolor{lmugreen}{\textbf{0.515}} & \textcolor{lmugreen}{0.268} \\
        LinearGAM (log)  & 0.444          & 0.506          & \textbf{0.255} \\
        GammaGAM         & 0.475          & 0.434          & 0.277 \\
        \bottomrule
      \end{tabular}\\[0.4em]
      \footnotesize Outer-fold means, 3-fold nested CV\@.\\
      \textcolor{lmugreen}{\textbf{LinearGAM selected.}}
      }
    \end{column}
  \end{columns}
\end{frame}

%--- Slide 14 ------------------------------------------------------------------
\begin{frame}{Final Model Performance (OOF)}
  \begin{columns}[t]
    \begin{column}{0.48\textwidth}
      \textbf{GAM (LinearGAM)}
      \begin{itemize}
        \item Fitted on raw \texttt{ml\_ade}
        \item Additive smooth terms, one per feature
      \end{itemize}
      \vspace{0.5em}
      \renewcommand{\arraystretch}{1.3}
      \begin{tabular}{lc}
        \toprule
        \textbf{Metric} & \textbf{OOF value} \\
        \midrule
        RMSE & 0.440 \\
        R$^2$ & 0.515 \\
        MAE  & 0.268 \\
        \bottomrule
      \end{tabular}
    \end{column}
    \begin{column}{0.48\textwidth}
      \textbf{XGBoost}
      \begin{itemize}
        \item Gradient-boosted trees
        \item Post-hoc interpretation via SHAP
      \end{itemize}
      \vspace{0.5em}
      \renewcommand{\arraystretch}{1.3}
      \begin{tabular}{lc}
        \toprule
        \textbf{Metric} & \textbf{OOF value} \\
        \midrule
        RMSE & \textbf{0.387} \\
        R$^2$ & \textbf{0.624} \\
        MAE  & \textbf{0.213} \\
        \bottomrule
      \end{tabular}
    \end{column}
  \end{columns}
  \vspace{0.5em}
  \footnotesize
  XGBoost outperforms GAM on all metrics. Both models are used as a
  cross-check: the feature-effect narratives should agree across approaches.
\end{frame}
